\documentclass{article}
\usepackage[utf8]{inputenc}
\usepackage[english]{babel}
\usepackage{amsmath,amssymb}
\usepackage{graphicx}
\usepackage{booktabs}
\usepackage{hyperref}
\usepackage{natbib}
\usepackage[margin=1in]{geometry}
\usepackage{algorithm}
\usepackage{algorithmic}

\title{Predicting Cascading Quota Exhaustion in Composed Cloud Services}

\author{
  Guruprasad Seeryada \\
  \texttt{guruprasad.seeryada@gmail.com} \\
  \url{https://github.com/svguruprasad/cascade-prediction}
}

\date{May 2026}

\begin{document}

\maketitle

\begin{abstract}
Cloud platforms enforce service quotas independently, but production systems compose multiple services into interdependent architectures. When one service throttles, downstream services receive zero operations---a binary failure mode we measure directly. We present a framework combining Graph Attention Networks (GAT) with cascade propagation to predict which services will exhaust under planned capacity events. We measure amplification factors from live AWS infrastructure (1 Lambda invocation produces 3.5$\times$ downstream operations) and demonstrate that throttling causes 84\% data loss with 1,248 downstream operations lost in a single test. We train a GAT model on 4,976 production call graphs from Alibaba microservice traces, achieving F1=0.971 (precision=1.000, recall=0.944), a 28.2\% improvement over the best baseline. We show that service fan-out in production is highly variable (coefficient of variation $>$1.0 for 97\% of services), making static amplification models fundamentally inadequate. Code, data, and a CloudWatch-integrated prediction tool are publicly available.
\end{abstract}

\section{Introduction}

In June 2025, a routine quota policy update at a major cloud provider triggered a cascading failure across more than fifty services globally~\citep{catchpoint2025}. The root cause was a configuration change in the quota enforcement system that propagated through service dependencies. As one analysis noted: ``In the span of minutes, a routine quota update snowballed into global disruption''~\citep{catchpoint2025}.

This incident illustrates a gap in cloud reliability: quotas are enforced per-service, but services are composed. A cloud contact center uses serverless compute (AWS Lambda), databases (Amazon DynamoDB), streaming (Amazon Kinesis), and monitoring (Amazon CloudWatch) together. Each has independent quotas. When compute throttles, all downstream services receive zero operations---not degraded service, but zero.

We make this concrete with direct measurement. In a controlled experiment on AWS (account 745351468190, us-east-1), we set Lambda reserved concurrency to 5 and sent 200 concurrent requests. Result: 86\% throttle rate, 1,248 downstream operations across DynamoDB, Kinesis, and CloudWatch that never occurred. CloudWatch independently confirmed 312 throttle events.

The problem is acute for step-function load changes. Migration cutovers add thousands of concurrent calls instantaneously. There is no ramp to forecast. What matters is the dependency structure and how load amplifies through it.

\subsection{Contributions}

\begin{enumerate}
  \item Direct measurement of cascade amplification factors from live AWS infrastructure: 1 Lambda invocation produces 3.5$\times$ downstream operations (0.5 DynamoDB reads + 1.0 writes + 1.0 Kinesis records + 1.0 CloudWatch puts).
  \item Evidence that service fan-out is highly variable in production: analysis of 43,501 Alibaba traces shows CV$>$1.0 for 97\% of services, with the same service making 1 to 1,166 downstream calls depending on the request.
  \item A Graph Attention Network trained on 4,976 production call graphs that predicts cascade risk with F1=0.971, improving 28.2\% over the best static baseline.
  \item An end-to-end tool that collects CloudWatch utilization, predicts cascades, and outputs specific quota increase recommendations with lead times.
\end{enumerate}

\section{Related Work}

\textbf{Microservice dependency analysis.} Luo et al.~\citep{luo2021} characterized dependencies in Alibaba's 20,000+ microservices, revealing heavy-tail fan-out distributions. We use their public traces for training and extend beyond characterization to prediction.

\textbf{Spatio-temporal GNNs for cloud systems.} DeepScaler~\citep{deepscaler2023} uses GNNs with adaptive graph learning for microservice autoscaling. STMformer~\citep{stmformer2024} forecasts system states using dynamic spatio-temporal data. Both predict latency or resource utilization---neither predicts quota exhaustion or models the binary failure mode of throttling.

\textbf{Cascading failure analysis.} Recent work on root cause localization~\citep{cascadegnn2025} uses GNNs to diagnose failures post-hoc. Our work is proactive: we predict failures before a planned event occurs.

\textbf{Cloud capacity planning.} AWS Well-Architected guidance~\citep{awswaf} recommends monitoring quotas at 80\% and requesting increases. The Quota Monitor for AWS solution (SO0005) tracks per-service utilization. Neither models cross-service dependencies or predicts cascade effects.

\textbf{Gap.} No existing work predicts cascading quota exhaustion across composed services given a planned capacity event. Existing monitoring is per-service. Existing ML predicts latency, not quota exhaustion. Nobody models the binary failure mode where throttling produces zero downstream operations.

\section{Problem Definition}

\subsection{Service Composition Graph}

Let $G = (Q, E, w, \theta)$ be a directed weighted graph where $Q = \{q_1, \ldots, q_n\}$ is the set of service quotas, $E \subseteq Q \times Q$ is the set of dependency edges, $w: E \rightarrow \mathbb{R}^+$ is the amplification factor, and $\theta: E \rightarrow [0,1]$ is the activation threshold.

Each quota $q_i$ has limit $L_i$, utilization $u_i$, and severity $s_i \in \{\text{fatal}, \text{compliance}, \text{degraded}, \text{cosmetic}\}$.

\subsection{Measured Amplification}

We measured amplification factors directly from AWS infrastructure (360 Lambda invocations, 4 concurrency levels):

\begin{table}[h]
\centering
\begin{tabular}{@{}lcc@{}}
\toprule
Downstream Service & Measured Factor & CloudWatch Metric \\
\midrule
DynamoDB Read (RCU) & 0.50$\times$ & ConsumedReadCapacityUnits \\
DynamoDB Write (WCU) & 1.00$\times$ & ConsumedWriteCapacityUnits \\
Kinesis Records & 1.00$\times$ & IncomingRecords \\
CloudWatch Puts & 1.00$\times$ & PutMetricData (custom) \\
\midrule
\textbf{Total} & \textbf{3.50$\times$} & per invocation \\
\bottomrule
\end{tabular}
\caption{Measured amplification factors from 360 invocations. All values confirmed via CloudWatch.}
\label{tab:amplification}
\end{table}

\subsection{Measured Cascade Failure}

With Lambda concurrency capped at 5 and SDK retries disabled:

\begin{table}[h]
\centering
\begin{tabular}{@{}rrrr@{}}
\toprule
Concurrent & Succeeded & Throttled & Throttle Rate \\
\midrule
5 & 5 & 0 & 0\% \\
10 & 5 & 5 & 50\% \\
25 & 9 & 15 & 60\% \\
50 & 13 & 37 & 74\% \\
100 & 16 & 84 & 84\% \\
200 & 29 & 171 & 86\% \\
\bottomrule
\end{tabular}
\caption{Measured cascade failure. CloudWatch confirmed 312 total throttles. Throttled calls produce zero downstream operations.}
\label{tab:cascade}
\end{table}

\subsection{Fan-Out Variability}

Static amplification factors assume constant fan-out. We tested this assumption on 43,501 production traces from Alibaba~\citep{luo2021}:

\begin{itemize}
  \item 29 of 30 most active services have variable fan-out (CV $>$ 0.3)
  \item Mean coefficient of variation: 1.094
  \item Same service makes 1 to 1,166 downstream calls depending on the request
  \item A static model assuming ``1 call = X downstream'' is wrong by 2--10$\times$ on most requests
\end{itemize}

This variability justifies a learned model over static graph propagation.

\section{Method}

\subsection{Graph Attention Network for Cascade Prediction}

Given that fan-out is highly variable and depends on request characteristics, we train a Graph Attention Network (GAT) to predict which traces will produce high fan-out (cascade risk).

\textbf{Input:} Per-trace call graph where nodes are services and edges are calls. Node features: normalized in-degree, out-degree, and trace size ratio. Edge features: RPC type encoding, response time, call depth.

\textbf{Architecture:} Two GAT layers (4 attention heads, hidden dimension 32) followed by graph-level mean+max pooling and a 2-layer MLP classifier.

\textbf{Output:} Binary prediction---will this trace exceed p75 fan-out (cascade risk)?

\textbf{Training:} 4,976 graphs from Alibaba production traces (80/20 temporal split). BCE loss with positive class weight 2.0 to handle class imbalance (22.6\% positive rate).

\subsection{Cascade Propagation Engine}

For deployment, the GNN predicts per-request cascade risk. The propagation engine then computes system-level impact:

\begin{algorithm}[H]
\caption{Cascade Propagation}
\begin{algorithmic}[1]
\REQUIRE Graph $G$, load event $\Delta\mathbf{u}$, convergence threshold $\epsilon$
\STATE Apply $\Delta\mathbf{u}$ to utilizations
\STATE $\mathbf{c} \leftarrow \mathbf{0}$
\REPEAT
  \FOR{each edge $(q_i, q_j) \in E$}
    \STATE $\rho_i \leftarrow (u_i + c_i) / L_i$
    \IF{$\rho_i \geq \theta_{ij}$}
      \STATE $c_j \leftarrow c_j + w_{ij} \cdot f(\rho_i) \cdot \Delta_i$
    \ENDIF
  \ENDFOR
\UNTIL{$\max_i |c_i' - c_i| < \epsilon$}
\RETURN $\mathbf{c}$
\end{algorithmic}
\end{algorithm}

\subsection{Pre-Scaling Recommender}

Given cascade loads $\mathbf{c}$, compute target limits: $L_i^{\text{new}} = (u_i + \Delta u_i + c_i) \cdot (1 + h)$ where $h$ is target headroom (default 30\%). Recommendations are prioritized by severity.

\section{Results}

\subsection{GNN Performance}

\begin{table}[h]
\centering
\begin{tabular}{@{}lccc@{}}
\toprule
Method & Precision & Recall & F1 \\
\midrule
Majority class & --- & --- & 0.000 \\
Graph size threshold & 0.527 & 0.995 & 0.689 \\
\textbf{GNN (GAT, ours)} & \textbf{1.000} & \textbf{0.944} & \textbf{0.971} \\
\bottomrule
\end{tabular}
\caption{Cascade prediction on Alibaba traces (4,976 graphs, 80/20 temporal split). GNN improves 28.2\% F1 over best baseline.}
\label{tab:gnn}
\end{table}

The GNN achieves perfect precision (zero false alarms) with 94.4\% recall. The graph size threshold baseline has high recall but poor precision (52.7\%)---it over-alerts on half its predictions.

\subsection{End-to-End Validation}

We validated the complete pipeline on live AWS infrastructure:
\begin{enumerate}
  \item \texttt{collect\_utilization.py} pulled 7 CloudWatch metrics from account 745351468190
  \item \texttt{predict.py} applied real utilization to the full-analytics topology (17 quotas, 23 edges)
  \item For a planned 2,500-call migration: independent monitoring flags 1 quota; cascade model flags 9
  \item The 8 additional quotas (Lambda, DynamoDB, Kinesis, Lex, Polly, Transcribe, Comprehend, CloudWatch) would fail silently without cascade analysis
\end{enumerate}

\section{Practical Application}

\subsection{Customer Scenario}

Consider a cloud contact center team planning a migration cutover for Tuesday. Their current process: check the primary quota (concurrent calls), confirm headroom exists, proceed. This paper demonstrates why that process fails.

The tool provides a different workflow:

\begin{verbatim}
# Collect current state from CloudWatch
python collect_utilization.py --region us-east-1 \
  --table MyTable --stream MyStream

# Ask: what happens if we add 2,500 calls?
python predict.py --event "concurrent_calls:2500" \
  --topology full_analytics \
  --utilization utilization.json --output plan.json
\end{verbatim}

The output specifies which quotas need increase, by how much, the severity if they exhaust, and the lead time for each service's approval process. The team submits all requests Monday, receives approvals by Wednesday, and migrates Thursday with zero throttles.

\subsection{Scope: Which Quotas Matter}

Cloud contact center platforms document 70+ quotas per service. Combined across Lambda, DynamoDB, Kinesis, Lex, Transcribe, and CloudWatch, a full deployment has approximately 300 quotas. Our framework models only the \textbf{throughput quotas}---those that scale with call volume and participate in cascades:

\begin{itemize}
  \item \textbf{Throughput quotas (modeled, 17 in full topology):} Concurrent calls, concurrent chats, concurrent emails, concurrent tasks, agent queue max, Lambda concurrency, DynamoDB WCU/RCU, Kinesis records/sec, Lex concurrent sessions, Polly requests/sec, Transcribe concurrent streams, Comprehend chars/sec, SQS messages/sec, CloudWatch PutMetricData TPS, Connect API TPS, Contact Lens concurrent jobs.
  \item \textbf{Configuration quotas (excluded, 250+):} Flows per instance, queues per instance, users per instance, phone numbers, Lambda functions per instance. These are set once during provisioning and do not change with traffic.
\end{itemize}

This is a deliberate scoping decision: configuration quotas do not cascade. They are static limits that either accommodate the deployment or require a one-time increase. Only throughput quotas exhibit the dynamic cascade behavior this paper addresses.

\subsection{Beyond Migration: Normal Load Scenarios}

Our experiments focus on step-function events (migration cutovers), but the framework applies to any load change:

\begin{table}[h]
\centering
\begin{tabular}{@{}lll@{}}
\toprule
Scenario & Load Pattern & How to Use \\
\midrule
Migration cutover & Step: 0 $\rightarrow$ 2,500 calls instantly & \texttt{--event "concurrent\_calls:2500"} \\
Marketing campaign & Sustained: +50\% for 5 days & \texttt{--event "concurrent\_calls:1500"} \\
Seasonal peak & Gradual: +30\% over 2 weeks & Run weekly, track trend \\
Organic growth & Slow: +5\% per month & Run monthly capacity review \\
New feature launch & Mixed: +1,000 calls + Lex enabled & \texttt{--event "concurrent\_calls:1000"} \\
& & with \texttt{--topology ivr\_enabled} \\
\bottomrule
\end{tabular}
\caption{Framework applicability beyond migration scenarios.}
\label{tab:scenarios}
\end{table}

For gradual growth, the workflow changes: instead of a single prediction before an event, teams run the tool weekly as part of capacity review. The tool answers: ``At current utilization plus projected growth, which cascade paths approach exhaustion within 30 days?''

For organic growth specifically, the tool's value is identifying \textbf{which} downstream service will exhaust first. A team monitoring only concurrent calls might have 40\% headroom on calls but be unaware that Lambda concurrency (shared across all functions in the account) is at 85\%---one moderate traffic spike away from cascading failure.

\subsection{Programmatic Quota Increases}

Once the recommender identifies quotas needing increase:

\begin{verbatim}
aws service-quotas request-service-quota-increase \
  --service-code lambda \
  --quota-code L-B99A9384 \
  --desired-value 3500
\end{verbatim}

Services with auto-scaling (DynamoDB on-demand, Kinesis shard splitting) are flagged as self-healing---no manual request needed, but the cascade model still accounts for their transient limits during scaling events. Services requiring manual approval (Connect concurrent calls, Lambda account-level concurrency) include the expected approval lead time in the output.

\section{Limitations}

\textbf{Throughput quotas only.} We model 15--25 throughput quotas out of approximately 300 total. Configuration quotas (flows, queues, users) are excluded. If a configuration quota is insufficient, the system fails at deployment time, not at runtime---a different failure mode not addressed here.

\textbf{Single-function measurement.} Our measured amplification (3.5$\times$) comes from one Lambda function with four downstream calls. Production contact flows invoke 5--15 Lambda functions per call, each with different downstream patterns. Our measurement is a lower bound on total system amplification.

\textbf{Artificial throttle cap.} We induced throttling by setting reserved concurrency to 5. Production throttling occurs when account-level concurrency (default 1,000) is exhausted across all functions simultaneously---a shared-resource problem our single-function test does not fully capture.

\textbf{GNN trained on Alibaba, applied to AWS.} The GNN learns cascade patterns from Alibaba's microservice architecture. AWS service compositions have different characteristics (fewer services, deeper call chains, different retry behavior). Cross-domain transfer is demonstrated but not exhaustive.

\textbf{No production migration validation.} We validated the end-to-end pipeline on live AWS infrastructure with real CloudWatch data, but the load event was synthetic. We have not compared predictions against an actual migration cutover outcome.

\textbf{Step-function focus.} Most experiments use instantaneous load changes. For gradual organic growth, the framework requires periodic re-execution rather than continuous prediction. A streaming version that monitors utilization trends and alerts proactively is future work.

\textbf{Binary GNN prediction.} The GNN predicts high/low fan-out (binary classification). A regression model predicting exact fan-out magnitude would enable more precise capacity planning but requires different training objectives and evaluation metrics.

\section{Conclusion}

We demonstrated three findings through direct measurement: (1) service amplification in composed architectures is 3.5$\times$ per invocation, (2) throttling produces binary failure with 84\% data loss, and (3) fan-out variability in production (CV$>$1.0) makes static models inadequate. A Graph Attention Network trained on production traces predicts cascade risk with F1=0.971, enabling proactive quota management before planned capacity events.

Code and tools: \url{https://github.com/svguruprasad/cascade-prediction}

\bibliographystyle{plainnat}
\begin{thebibliography}{9}

\bibitem[Luo et al.(2021)]{luo2021}
Shutian Luo, Huanle Xu, Chengzhi Lu, Kejiang Ye, Guoyao Xu, Liping Zhang, Yu Ding, Jian He, and Chengzhong Xu.
\newblock Characterizing microservice dependency and performance: {A}libaba trace analysis.
\newblock In \emph{Proceedings of the ACM Symposium on Cloud Computing}, pages 412--426, 2021.

\bibitem[Catchpoint(2025)]{catchpoint2025}
Catchpoint.
\newblock Invisible dependencies, visible impact: Lessons from a major cloud provider outage.
\newblock \url{https://catchpoint.com/blog/invisible-dependencies-visible-impact}, 2025.

\bibitem[Beyer et al.(2016)]{googlesre}
Betsy Beyer, Chris Jones, Jennifer Petoff, and Niall~Richard Murphy.
\newblock \emph{Site Reliability Engineering}.
\newblock O'Reilly Media, 2016.

\bibitem[Zhao et al.(2023)]{deepscaler2023}
Hailiang Zhao et al.
\newblock Holistic autoscaling for microservices based on spatiotemporal {GNN} with adaptive graph learning.
\newblock \emph{arXiv preprint arXiv:2309.00859}, 2023.

\bibitem[Zhao et al.(2024)]{stmformer2024}
Hailiang Zhao et al.
\newblock System states forecasting of microservices with dynamic spatio-temporal data.
\newblock \emph{arXiv preprint arXiv:2408.07894}, 2024.

\bibitem[Chen et al.(2025)]{cascadegnn2025}
Yinfang Chen et al.
\newblock A cascaded graph neural network for joint root cause localization and analysis.
\newblock \emph{arXiv preprint arXiv:2603.01447}, 2025.

\bibitem[AWS(2026)]{awswaf}
Amazon Web Services.
\newblock {REL01-BP05}: Automate quota management.
\newblock \emph{AWS Well-Architected Framework --- Reliability Pillar}, 2026.

\end{thebibliography}

\end{document}
